% !TEX root = ../main.tex
\chapter{Tangent Spaces}\label{ch:c05}
\begin{paracol}{2}

\begin{sources}
\begin{tabularx}{\linewidth}{@{}l l L@{}}
\toprule
\textbf{Source} & \textbf{Length} & \textbf{Content}\\
\midrule
\seg{5}{1}{Lecture 5: Tangent Spaces} & 1:36:41 & The velocity of a curve as a linear map on smooth functions; the tangent space as the set of all velocities, intrinsic versus embedded pictures; the vector space structure (scalar multiples by reparametrising, sums through a chart) with student questions;
the dimension of velocity components; components of a velocity and the chart-induced basis; change of components under a change of chart and a remark on special relativity; the cotangent space, the gradient $\mathrm df$, the dual basis $\mathrm dx^i$ and the transformation of covector components.\\
\seg{5}{2}{Tutorial 5: Tangent Spaces} & 1:32:26 & Solutions of the sheet by a tutor (not Schuller): fifteen true/false statements; the symbol $\partial/\partial x^i|_p$ and its chain rule, which makes $\partial f/\partial x^i|_p$ a $(0,1)$-tensor; components of a velocity on $\R^2$ with two charts, computed directly and by the transformation law;
the gradient annihilates velocities of curves in a level set; why a sum of curves is not well defined, but the sum of velocities is.\\
\bottomrule
\end{tabularx}

\smallskip
\textbf{Transcripts} \file{materials/transcripts/C05-V01.txt} and \file{C05-V02.txt} (YouTube captions).
\textbf{Sheet 5:} see \url{https://tales.mbivert.com/on-heraeus-winter-school-tutorials/}.
\end{sources}

\begin{reading}[Where this chapter sits in the other notes]
\emph{Geometric Anatomy of Theoretical Physics} notes: ``Construction of the tangent bundle'' is chapter 10 there, and the tangent spaces themselves (velocities of curves as derivations, the chart-induced basis,
the change of components, the cotangent space) are chapter 9: ``Differentiable structures, tangent spaces'' [title not checked]. Lecture 6 of this book (chapter~\ref{ch:c06}) builds the tangent bundle and the vector fields on it, i.e.\ chapters 10 and 11 there. The vector spaces, duals
and tensors used throughout are chapter~\ref{ch:c03} here; the smooth manifolds and chart representatives are chapter~\ref{ch:c04}.
\end{reading}

% ==================================================================
\section{The question: what is the velocity of a curve?}\label{sec:c05-velocity}

Welcome to the fifth lecture \ts{5}{1}{0:00}. Under this unassuming name ``goes really a lot of the basic concepts, mathematically, that underlie general relativity'', so this is a really important lecture. The lead question: on a manifold, if one considers a curve, what is the \emph{velocity} of that curve at the point $p$?
Velocity and not speed: speed cannot be defined before a so-called metric is defined on the manifold, which comes later. One starts with velocity, and it is tradition to call the velocity a vector, and in the formalism it will indeed be a vector; but ``I ask you at first to completely erase anything you ever learned about velocities, we are trying to rediscover it''
\ts{5}{1}{0:20}\ts{5}{1}{0:41}\ts{5}{1}{1:02}. Rather than making motivating remarks that would confuse, one starts with the definition and then discusses it \ts{5}{1}{1:22}.

\begin{definition}[New concepts: velocity of a curve at a point, $C^\infty(M)$]\label{def:c05-velocity}
Let $(M,\mathcal O,\mathcal A)$ be a smooth manifold and $\gamma\colon\R\to M$ a curve which is at least $C^1$ (once continuously differentiable with respect to the charts of the smooth atlas), and suppose that $\gamma(\lambda_0)=p$ for some parameter value $\lambda_0$ (otherwise the question for the velocity at $p$ is meaningless).
\begin{enumerate}
  \item $C^\infty(M)$ denotes the set of all smooth functions $f\colon M\to\R$, with the \emph{pointwise} operations
        \[
          (f+g)(p):=f(p)+g(p),\qquad(\lambda g)(p):=\lambda\,g(p),
        \]
        where the $+$ and the product on the right are those of $\R$. With these two operations $C^\infty(M)$ is a real vector space.
  \item The \textbf{velocity of $\gamma$ at $p$} is the linear map
        \[
          v_{\gamma,p}\colon C^\infty(M)\xrightarrow{\ \sim\ }\R,\qquad v_{\gamma,p}(f):=(f\circ\gamma)'(\lambda_0),
        \]
        the derivative of the real function $f\circ\gamma\colon\R\to\R$ at the parameter value $\lambda_0$.
\end{enumerate}
\ts{5}{1}{1:44}\ts{5}{1}{2:26}\ts{5}{1}{3:07}\ts{5}{1}{3:52}\ts{5}{1}{4:34}\ts{5}{1}{4:55}.
\end{definition}
So velocities are \emph{linear maps} (``something you need to get used to''): $v_{\gamma,p}$ eats an arbitrary smooth function and returns a real number, namely how fast the function changes along the curve. The linear map goes from one vector space, $C^\infty(M)$, to another, $\R$ (a vector space over itself with its own addition and multiplication)
\ts{5}{1}{2:46}\ts{5}{1}{4:13}.

\paragraph{Intuition: the child that runs around.} This is ``exactly what children, uncorrupted by high school physics, capture as the concept of velocity'': a child runs through the world and observes some temperature function $f$ (level lines drawn on the board); $f\circ\gamma$ is just a function from time to temperature, and
what one feels is the rate of change of the temperature as one runs along the curve. So the velocity of a curve at $p$ is the \emph{directional derivative} along the curve, for any $f$ \ts{5}{1}{5:15}\ts{5}{1}{5:36}\ts{5}{1}{5:58}\ts{5}{1}{6:19}\ts{5}{1}{6:39}.

Compare with the old picture: one takes the gradient of $f$ and dots the vector into it, $\vec v\cdot\nabla f$ (the vector being a list of components); or one says that $v^i\partial_i$ is the differential operator acting on $f$. The only shift of philosophy is that rather than saying the \emph{components} are the vector, in differential geometry one says that
the \emph{map} is the vector: a linear map that takes an $f$ and maps it to a number. Vectors survive as the directional derivatives they induce, and if one knows how the map acts on any function one can make it act on the coordinate functions, get Kronecker deltas, and recover all the components. ``Now delete everything I said, because everything I said will now be made precise''
\ts{5}{1}{7:00}\ts{5}{1}{7:21}\ts{5}{1}{7:42}\ts{5}{1}{8:02}\ts{5}{1}{8:23}\ts{5}{1}{8:43}.
\begin{coursenote}[Derivations]
One checks at once that $v_{\gamma,p}(fg)=f(p)\,v_{\gamma,p}(g)+g(p)\,v_{\gamma,p}(f)$ (product rule for $(fg)\circ\gamma$): velocities are \emph{derivations} at $p$, which is how other texts define tangent vectors. Schuller does not mention this in the lecture. (Not said in class.)
\end{coursenote}

\section{The tangent space}\label{sec:c05-tangent}

\begin{definition}[New concepts: tangent space $T_pM$, tangent vector]\label{def:c05-tangent}
For each point $p$ of a smooth manifold $M$ the \textbf{tangent space to $M$ at $p$} is the set
\[
  T_pM:=\{\,v_{\gamma,p}\mid\gamma\text{ a smooth curve through }p\,\},
\]
the collection of all velocities of all smooth curves through $p$. Its elements are called \textbf{tangent vectors} at $p$
\ts{5}{1}{9:03}\ts{5}{1}{9:24}\ts{5}{1}{9:48}.
\end{definition}
The picture: a point and the curves through it, each one with a tangent vector (a directional derivative) attached. A curve through the same point running at double parameter speed gives a different directional derivative, because it scales; a curve that is different from another one may nevertheless give the same tangent vector \ts{5}{1}{10:10}\ts{5}{1}{10:31}.
(The name: ``tangent'' from the Latin \emph{tangere}, to touch [?, not said in class].)

\paragraph{Intrinsic versus embedded.} The definition makes \emph{no reference to an ambient space}. One may think of the tangent space as an arrow-plane touching the surface, as alluded to earlier, but that makes reference to an ambient space; the intrinsic picture is the cleverer one. In the beginning of differential geometry there was a debate for many years whether it is better to define objects
intrinsically, lying within the manifold, or to take the manifold embedded in a bigger space (the extrinsic point of view), and whether they lead to the same results; famous embedding theorems at all structural levels say that it does not matter, one always finds a big enough ambient space. So the debate is settled. Why prefer the intrinsic perspective?
If the manifold is finally the universe, all of spacetime, ``it would be a little bit uneconomical to refer to something around the universe''; it is better to define all concepts from within the universe, and technically it also points to the right ideas. ``Rather this than this as a picture, although really everybody thinks of it as the other as well; but technically you must not''
\ts{5}{1}{10:51}\ts{5}{1}{11:13}\ts{5}{1}{11:33}\ts{5}{1}{11:54}\ts{5}{1}{12:14}\ts{5}{1}{12:35}.

\section{The tangent space is a vector space}\label{sec:c05-vs}

So far $T_pM$ is just a set. Do the velocities at a point form a vector space? Yes, and it is a key observation. Because this is a technical key point that will be used all the time, Schuller shows it in detail: ``already in this calculation you learn all the techniques that you need in the future for your own calculations, so it is not just a proof that you then never look at again''
\ts{5}{1}{12:56}\ts{5}{1}{13:16}.

\begin{theorem}[$T_pM$ is a vector space]\label{thm:c05-vs}
Define, for velocities at the same point $p$ and $\alpha\in\R$,
\[
  (v_{\gamma,p}\oplus v_{\delta,p})(f):=v_{\gamma,p}(f)+v_{\delta,p}(f),\qquad(\alpha\odot v_{\gamma,p})(f):=\alpha\,v_{\gamma,p}(f)\qquad(f\in C^\infty(M)).
\]
Then both are again elements of $T_pM$, hence $(T_pM,\oplus,\odot)$ is a real vector space.
\end{theorem}
\ts{5}{1}{13:37}\ts{5}{1}{13:58}\ts{5}{1}{14:19}\ts{5}{1}{14:40}\ts{5}{1}{15:00}. Note that one cannot add a vector at $p$ to a vector at another point $q$; only two vectors \emph{at the same point}.
\begin{proof}
Because the elements are linear maps with values in $\R$, one can immediately define $\oplus$ and $\odot$ as above (the $+$ on the right is that of $\R$, the left one is the addition to be defined in $T_pM$). The sum and the multiple are then linear maps $C^\infty(M)\to\R$. But for these to be a vector space addition and scalar multiplication
they must \emph{close within} $T_pM$: the result has to lie in the specific subset of maps that arise from a curve. Hence it remains to show:
\begin{enumerate}
  \item for $v_{\gamma,p}$ and $\alpha\in\R$ there is a curve $\tau$ with $\alpha\odot v_{\gamma,p}=v_{\tau,p}$;
  \item for $v_{\gamma,p}$ and $v_{\delta,p}$ there is a curve $\sigma$ with $v_{\gamma,p}\oplus v_{\delta,p}=v_{\sigma,p}$.
\end{enumerate}
\ts{5}{1}{15:21}\ts{5}{1}{15:41}\ts{5}{1}{16:02}\ts{5}{1}{16:22}\ts{5}{1}{16:42}\ts{5}{1}{17:03}\ts{5}{1}{17:24}\ts{5}{1}{17:45}.

\emph{(i) Scalar multiples.} Let $\gamma(\lambda_0)=p$. Define $\tau:=\gamma\circ\mu_\alpha$ with $\mu_\alpha\colon\R\to\R$, $\mu_\alpha(\lambda):=\alpha\lambda+\lambda_0$ (one runs through the old curve at $\alpha$ times the parameter speed; the parameter is shifted so that $\tau(0)=\gamma(\lambda_0)=p$; Schuller first uses $\lambda$ also for the
multiplier, then switches to $\alpha$ and corrects the shift on the board). By the chain rule
\[
  \begin{aligned}v_{\tau,p}(f)&=(f\circ\gamma\circ\mu_\alpha)'(0)=(f\circ\gamma)'(\mu_\alpha(0))\cdot\mu_\alpha'(0)\&=(f\circ\gamma)'(\lambda_0)\cdot\alpha=\alpha\,v_{\gamma,p}(f).\end{aligned}
\]
So $\tau$ does the trick, and the construction needs no chart \ts{5}{1}{18:06}\ts{5}{1}{18:27}\ts{5}{1}{18:49}\ts{5}{1}{19:10}\ts{5}{1}{19:31}\ts{5}{1}{19:52}\ts{5}{1}{20:13}\ts{5}{1}{20:55}\ts{5}{1}{21:17}\ts{5}{1}{21:38}\ts{5}{1}{21:58}\ts{5}{1}{22:18}\ts{5}{1}{22:39}.

\emph{(ii) Sums.} Let $\gamma(\lambda_0)=p=\delta(\lambda_1)$ (the two curves may hit $p$ at different parameter values). Here one has to do ``something extremely ugly: make a choice of chart''. The moment one says this, the ill-definedness alarm bells have to ring: if in the end anything depends on the chart, ``I am not going to buy it'';
one is not against choosing charts, only against choosing charts and not knowing which assumption one is making; and here it turns out to be no assumption at all. Choose any chart $(U,x)$ with $p\in U$ (possible by the definition of a topological manifold) and define
\[\begin{aligned}
\sigma_x(\lambda)&:=x^{-1}\Bigl((x\circ\gamma)(\lambda_0+\lambda)\\
&\quad+(x\circ\delta)(\lambda_1+\lambda)-(x\circ\gamma)(\lambda_0)\Bigr),
\end{aligned}\]
where the sum takes place in $\R^d$ (the last term is a constant element of $\R^d$ that does not depend on $\lambda$, but is needed). The subscript $x$ is ``a constant reminder that this construction depends on the chart map''. First, $\sigma_x(0)=x^{-1}\bigl(x(\gamma(\lambda_0))+x(\delta(\lambda_1))-x(\gamma(\lambda_0))\bigr)=\delta(\lambda_1)=p$,
so the curve runs through $p$ at parameter value $0$. Next, for $f\in C^\infty(M)$ insert the identity $x^{-1}\circ x$:
\[\begin{aligned}
v_{\sigma_x,p}(f)&=(f\circ\sigma_x)'(0)\\
&=\bigl((f\circ x^{-1})\circ(x\circ\sigma_x)\bigr)'(0).
\end{aligned}\]
Here $f\circ x^{-1}$ goes from a piece of $\R^d$ to $\R$ and $x\circ\sigma_x$ from $\R$ to $\R^d$: this is undergraduate analysis, and the multi-dimensional chain rule gives (summation over $i$)
\[\begin{aligned}
v_{\sigma_x,p}(f)&=(x\circ\sigma_x)^{i\,\prime}(0)\\
&\quad\cdot\partial_i(f\circ x^{-1})\bigl(x(p)\bigr).
\end{aligned}\]
Now $x\circ\sigma_x$ is the sum in the chart, so the constant drops out under differentiation and
\[
  (x\circ\sigma_x)^{i\,\prime}(0)=(x\circ\gamma)^{i\,\prime}(\lambda_0)+(x\circ\delta)^{i\,\prime}(\lambda_1).
\]
Distributing the factor $\partial_i(f\circ x^{-1})(x(p))$ over the two terms and using the chain rule backwards on each,
\[\begin{aligned}
v_{\sigma_x,p}(f)&=(f\circ\gamma)'(\lambda_0)+(f\circ\delta)'(\lambda_1)\\
&=v_{\gamma,p}(f)+v_{\delta,p}(f).
\end{aligned}\]
No assumption on $f$ was made, so $v_{\sigma_x,p}=v_{\gamma,p}\oplus v_{\delta,p}$; and \emph{although a chart was chosen intermediately, neither side depends on it}. (The other vector space axioms are inherited from $\R$: addition and multiplication of maps with values in $\R$ are defined pointwise, and one uses
commutativity, associativity, the distributive laws and the unit of $\R$; see the student question below.)
\ts{5}{1}{23:00}\ts{5}{1}{23:23}\ts{5}{1}{23:43}\ts{5}{1}{24:04}\ts{5}{1}{24:24}\ts{5}{1}{24:45}\ts{5}{1}{25:06}\ts{5}{1}{25:27}\ts{5}{1}{25:48}\ts{5}{1}{26:08}\ts{5}{1}{26:30}\ts{5}{1}{26:50}\ts{5}{1}{27:12}\ts{5}{1}{27:32}\ts{5}{1}{27:52}\ts{5}{1}{28:14}\ts{5}{1}{28:35}\ts{5}{1}{28:56}\ts{5}{1}{29:16}\ts{5}{1}{29:37}\ts{5}{1}{29:58}\ts{5}{1}{30:18}\ts{5}{1}{30:38}\ts{5}{1}{31:01}\ts{5}{1}{31:22}\ts{5}{1}{31:42}\ts{5}{1}{32:02}\ts{5}{1}{32:23}\ts{5}{1}{32:43}\ts{5}{1}{33:04}\ts{5}{1}{33:25}\ts{5}{1}{33:46}.
\end{proof}

\paragraph{Student question: do I really not need the existence of a special chart?} A student did not get why the existence of a curve has been proven, and recalls that someone pointed out that such a chart might not exist. Schuller: what is needed is a candidate curve, constructed by using a chart; any chart that covers $p$ will do, and one exists by the very definition of a topological manifold.
The only requirement is that the \emph{result} be chart independent, and it is: with another chart $y$ everything would be written with $y$ and work out the same \ts{5}{1}{34:07}\ts{5}{1}{34:28}\ts{5}{1}{34:49}\ts{5}{1}{35:10}.

\paragraph{Student question: why the intrinsic picture, if charts map to Euclidean space?} A student: one says it is better to think intrinsically, but here one maps the manifold to real space, so locally one describes everything by a Euclidean space. Schuller: the idea is right (``I would not agree with the Euclidean'', the flat structure is the vector space structure of $\R^d$, which one only uses as a coordinate space);
a topological manifold \emph{locally} maps continuously and one-to-one to some $\R^d$. By the embedding picture he meant something else: one chooses an $\R^{1000}$ to look at a five-dimensional manifold, or an $\R^7$ to embed a three-dimensional one, so that what the intrinsic construction gives with only $\R^3$'s is what one would get as tangent spaces by working in the $\R^7$. That is ``even worse'': it really goes to an ambient space around the manifold, whereas the manifold definition only requires that pieces of the manifold go to the $\R^d$ of the dimension of the manifold.
And an arbitrary topological space has no charts in the first place: ``we make a huge assumption, we say the whole universe, locally, looks like $\R^d$, a massive assumption, and that is the assumption that spacetime is a topological manifold'' \ts{5}{1}{35:30}\ts{5}{1}{35:51}\ts{5}{1}{36:12}\ts{5}{1}{36:33}\ts{5}{1}{36:54}\ts{5}{1}{37:15}.

\paragraph{Student question: where did $x$ disappear in the first line?} Schuller inserts one more line of the chain rule backwards, $f\circ x^{-1}\circ x\circ\gamma$, in which the $x^{-1}\circ x$ vanishes again, ``very importantly'' \ts{5}{1}{37:36}\ts{5}{1}{37:57}.

\paragraph{Pictures of the two proofs.} A student was confused by the first proof (why no chart there). The picture: in the real world a curve with equal intervals of the parameter marked; one may think of the velocity as a little arrow, ``a bad picture, but not so bad for understanding the proof'': to draw \emph{double} the velocity one runs through the same curve at double parameter speed, which is what $\mu_2$ does,
and the tangent to $\gamma\circ\mu_2$ is twice the tangent to $\gamma$; no chart is needed (one could do it in a chart, but that would not be economical). The \emph{sum} is a real beast: the two curves must first be pushed down to a chart, where one defines the sum curve $x\circ\sigma_x$ and pulls it back by $x^{-1}$. With another chart $y$ the same two curves
give a different sum curve $\sigma_y$ through the same point. The curves $\sigma_x$ and $\sigma_y$ \emph{differ}, but they have the same tangent at $p$: ``the sum of curves is bad, but its derivative at the point of intersection is good'' (and one only uses it at this point). This is why one must not add curves: one could then add the trajectories of particles in a chart, which would amount to adding position vectors, which do not exist.
One cannot add curves on the manifold directly either: for $\gamma(\lambda)$ and $\delta(\lambda)$ landing on $M$ there is no ``$+$'' on the target. ``Positions do not constitute a vector space, it is always the same thing''
\ts{5}{1}{38:18}\ts{5}{1}{38:38}\ts{5}{1}{38:59}\ts{5}{1}{39:19}\ts{5}{1}{39:39}\ts{5}{1}{39:59}\ts{5}{1}{40:20}\ts{5}{1}{40:41}\ts{5}{1}{41:03}\ts{5}{1}{41:23}\ts{5}{1}{41:44}\ts{5}{1}{42:05}\ts{5}{1}{42:25}\ts{5}{1}{42:46}\ts{5}{1}{43:07}\ts{5}{1}{43:29}.

\paragraph{Student question: why the constants and different arguments in $\sigma_x$?} The smart move in $\sigma_x$ is the use of different arguments $\lambda_0+\lambda$ and $\lambda_1+\lambda$ and the subtracted constant, so that in the end one has the two derivatives, each evaluated at its own parameter value, independent of $\lambda$. A student asks whether, since the whole expression is under $x^{-1}$, one is not chart dependent from the start, and could not pull the $x$ out already
in the definition. No: the $x^{-1}$ is \emph{not a linear map}, so the sum cannot be pulled apart; the $+$ is a $+$ ``in the context of'' $\R^d$ (``context-sensitive plus''). (Schuller repeats the question for the room.) \ts{5}{1}{43:50}\ts{5}{1}{44:10}\ts{5}{1}{44:31}\ts{5}{1}{44:51}\ts{5}{1}{45:12}.

\paragraph{Student question: do you not need to check all the other vector space axioms?} Yes, in principle, and one can do it oneself; it is always the same construction: since one adds maps whose target is $\R$, one inherits the addition from $\R$, and almost everything comes immediately from the axioms of the field $\R$ (the addition, the multiplication of $\R\setminus\{0\}$, the distributive law); with them one inherits all the vector space axioms in a few steps. ``Anything I do not show you has my quality guarantee: it is true''
\ts{5}{1}{45:53}\ts{5}{1}{46:14}\ts{5}{1}{46:35}.

\section{The dimension of velocity components}\label{sec:c05-dimension}

\begin{aside}[Components of a velocity have the dimension 1/second]
Schuller wants ``to provoke'' at this point. The velocity as defined is the mathematical object representing the physical velocity of a thing in the real world, fully independent of any coordinate choice and of a choice of length. The physical dimension of speed is metres per second, all agree; what is the physical dimension of the \emph{components} of a velocity vector? ``Also metres per second, isn't it?'' No: it is hertz, one over second
(the components are rates of change of coordinate functions, and coordinates need not be lengths). ``You will learn, and in the end you will say: oh dear, it is really true'', but only after the lecture of next Wednesday, where one will see where the missing metre comes from. It is an important point and it also stalled Einstein for many years, not having understood this first \ts{5}{1}{46:56}\ts{5}{1}{47:16}\ts{5}{1}{47:36}\ts{5}{1}{47:57}\ts{5}{1}{48:17}.
\end{aside}
\begin{coursenote}[Which lecture?]
``Next Wednesday'' is presumably the lecture on metric manifolds (speed needs a metric, as said at the start); that it is the place where the metre comes back is Schuller's announcement and has not been checked against that lecture [?].
\end{coursenote}

\section{Components with respect to a chart; the chart-induced basis}\label{sec:c05-components}

Now one reconciles the abstract notion of velocity with components. Components of a vector need a basis, so one asks what a basis of the tangent space is. In this lecture one picks a point $p$ and stays within one tangent space; the next lecture looks at all tangent spaces together and considers vector fields
\ts{5}{1}{48:17}\ts{5}{1}{48:39}\ts{5}{1}{49:00}.

\paragraph{The calculation.} Let $(U,x)$ be a chart of the smooth atlas and $\gamma\colon\R\to U$ the portion of a curve lying in $U$ with $\gamma(0)=p$ (for simplicity $\lambda_0=0$). To understand $v_{\gamma,p}$, let it act on an arbitrary $f$ and insert the chart, as before:
\[
  v_{\gamma,p}(f)=(f\circ\gamma)'(0)=\bigl((f\circ x^{-1})\circ(x\circ\gamma)\bigr)'(0)=(x^i\circ\gamma)'(0)\cdot\partial_i(f\circ x^{-1})\bigl(x(p)\bigr),
\]
where $x^i\circ\gamma=(x\circ\gamma)^i$ is the $i$-th component of the chart representative of the curve: the inner derivative first, then the outer derivative with respect to the $i$-th \emph{entry}, summed over $i$ \ts{5}{1}{49:20}\ts{5}{1}{49:40}\ts{5}{1}{50:01}\ts{5}{1}{50:22}\ts{5}{1}{50:43}\ts{5}{1}{51:04}.

\paragraph{A remark on notation for partial derivatives.} In undergraduate analysis a function $f\colon\R\to\R$ has a derivative $f'$, again a function, which one can evaluate at $x$ or at $y$; the derivative is called $f'$, and it would be silly to write $\mathrm df/\mathrm dx$ only because $f$ has been given by a rule in terms of $x$, as the abstract notion of derivative has nothing to do with the name of the variable.
So Schuller does \emph{not} write $\partial/\partial x$ for ordinary partial derivatives of maps $\R^d\to\R$, but $\partial_i$: ``all you need to know is with respect to which entry'' \ts{5}{1}{51:25}\ts{5}{1}{51:46}\ts{5}{1}{52:06}\ts{5}{1}{52:27}.

\begin{definition}[New concepts: the symbol $\partial/\partial x^i|_p$, components of a velocity]\label{def:c05-chartbasis}
Let $(U,x)$ be a chart in the smooth atlas, $p\in U$ and $f\in C^\infty(U)$ [smooth function defined near $p$]. Define
\[
  \frac{\partial f}{\partial x^i}\Big|_p:=\partial_i\bigl(f\circ x^{-1}\bigr)\bigl(x(p)\bigr),\qquad\text{so that}\qquad\frac{\partial}{\partial x^i}\Big|_p\colon f\longmapsto\partial_i(f\circ x^{-1})(x(p)).
\]
(The $x$ in the symbol remembers \emph{which} chart is used, to distinguish it from a $y$; the index $i$ ``down'' in the denominator is only notation and has nothing to do with the name of the entries; the $p$ is the point of evaluation.)
The numbers
\[
  \dot\gamma^i_x(0):=(x^i\circ\gamma)'(0)
\]
(the derivative of the $i$-th component of the chart representative of $\gamma$) are the \textbf{components of $v_{\gamma,p}$ with respect to the chart $x$}. Then the calculation above reads
\[
  v_{\gamma,p}(f)=\dot\gamma^i_x(0)\,\frac{\partial f}{\partial x^i}\Big|_p,\qquad\text{that is}\qquad v_{\gamma,p}=\dot\gamma^i_x(0)\,\frac{\partial}{\partial x^i}\Big|_p
\]
as maps on $C^\infty(M)$ [with $f$ restricted to $U$]
\ts{5}{1}{52:47}\ts{5}{1}{53:08}\ts{5}{1}{53:28}\ts{5}{1}{53:49}\ts{5}{1}{54:09}\ts{5}{1}{54:29}\ts{5}{1}{54:51}\ts{5}{1}{55:11}\ts{5}{1}{55:32}\ts{5}{1}{55:53}.
\end{definition}
Schuller stresses: the symbol looks like a partial derivative (it is designed to look like one, ``it makes it accessible to the masses'', and one can prove that it satisfies all the rules one wants from it), but \emph{the function $f$ here is a function on the manifold} (on $U$) and not on the chart; ``whenever you see this symbol, do not be deceived by its simplicity: you immediately need to translate it'' into $\partial_i(f\circ x^{-1})(x(p))$,
there is no alternative \ts{5}{1}{53:49}\ts{5}{1}{54:09}\ts{5}{1}{54:29}\ts{5}{1}{54:51}. It is tradition to delete the brackets and say ``a vector acts on whatever is behind it''; since the equation holds for all $f\in C^\infty(M)$ one may drop the $f$. Under the use of a chart the vector is its components (the derivatives of the chart representative of the curve) times basis elements:
the numbers $\dot\gamma^i_x(0)$ are the components and so the maps $\partial/\partial x^i|_p$ must be the basis elements of the space in which the vector lives, $T_pM$ \ts{5}{1}{56:14}\ts{5}{1}{56:34}\ts{5}{1}{56:54}\ts{5}{1}{57:15}\ts{5}{1}{57:35}.

\paragraph{The picture: polar coordinates and the bird.} The choice of a chart gives a choice of bases in the tangent spaces, which is ``fully irrelevant'': another chart would give another basis. Picture: vectors in the tangent spaces of the real world; map the points down by $x$, and pull the coordinate net back by $x^{-1}$.
In polar coordinates (with $x^1=r$, $x^2=\varphi$, like the generalised coordinates of classical mechanics, apart from some slits) the vectors of a circular motion have components $(0,1),(0,1),\dots$ in the chart, ``it looks dramatically different from the world, and depends on the chart''. A bird flying in a circle is an objective reality;
in a polar chart it runs along a \emph{straight} line, which it does not in reality. So the components of vectors depend heavily on the chart, and drawing the vectors in the chart (or ``up here'', in the world, with the chart's components) is a bad idea: one should only draw them in the chart; there is a velocity that the bird has, and in the chart it looks different \ts{5}{1}{57:35}\ts{5}{1}{58:21}\ts{5}{1}{58:44}\ts{5}{1}{59:05}\ts{5}{1}{59:26}\ts{5}{1}{59:46}\ts{5}{1}{1:00:09}\ts{5}{1}{1:00:29}\ts{5}{1}{1:00:50}\ts{5}{1}{1:01:10}\ts{5}{1}{1:01:31}.

\begin{definition}[New concepts: chart-induced basis]\label{def:c05-inducedbasis}
If $(U,x)$ is a chart in the smooth atlas, the elements $\dfrac{\partial}{\partial x^1}\Big|_p,\dots,\dfrac{\partial}{\partial x^d}\Big|_p$ of $T_pM$ are called the \textbf{chart-induced basis} of $T_pM$ (for $p\in U$).
\end{definition}
\ts{5}{1}{1:01:54}\ts{5}{1}{1:02:15}. The calculation above already showed that they \emph{generate} $T_pM$ (every vector is represented as a linear combination of them). The name ``basis'' is therefore a theorem:
\begin{theorem}[The chart-induced basis is a basis]\label{thm:c05-basis}
Let $(U,x)$ be a chart and $p\in U$. The vectors $\partial/\partial x^i|_p$, $i=1,\dots,d$, are linearly independent and span $T_pM$, hence form a basis of $T_pM$. In particular $\dim T_pM=d=\dim M$.
\end{theorem}
\begin{proof}
Spanning was shown above. For linear independence let $\lambda^i\,\partial/\partial x^i|_p=0$ (the zero map on functions). Apply this to the $j$-th coordinate function $x^j\colon U\to\R$, which is a smooth function with respect to the atlas (one learns in the tutorials: seen in the chart, $x^j\circ x^{-1}$ is a component of the identity). By the definition of the symbol
\[
  \frac{\partial x^j}{\partial x^i}\Big|_p=\partial_i\bigl(x^j\circ x^{-1}\bigr)\bigl(x(p)\bigr),
\]
and $x^j\circ x^{-1}\colon\R^d\to\R$ takes $(\alpha^1,\dots,\alpha^d)$ to $\alpha^j$, so its $i$-th partial derivative is $\delta^j_i$. Hence $0=\lambda^i\delta^j_i=\lambda^j$ for all $j$. Therefore there are $d$ basis vectors, and the vector space dimension of $T_pM$ is $d$.
\ts{5}{1}{1:02:38}\ts{5}{1}{1:02:58}\ts{5}{1}{1:03:19}\ts{5}{1}{1:03:40}\ts{5}{1}{1:04:01}\ts{5}{1}{1:04:22}\ts{5}{1}{1:05:45}\ts{5}{1}{1:06:06}\ts{5}{1}{1:06:28}\ts{5}{1}{1:06:50}\ts{5}{1}{1:07:10}\ts{5}{1}{1:07:34}\ts{5}{1}{1:07:55}.
\end{proof}
Two different notions of dimension meet here: the vector space dimension of $T_pM$ and the (topological) dimension of the manifold, which is a very different concept, and they agree: ``the tangent space at each point has the dimension of the manifold''. The $d$ came from the fact that there are $d$ coordinate maps.

\paragraph{Student question: $T_pM$ or $T_pU$?} A student asks why $T_pM$ and $T_pU$ differ. Schuller: $U$ is just a domain, an area, in which the chart is valid; the basis only exists there, and in general one cannot find a global chart for a manifold, and even when one can cover it by global charts one might not find a global coordinate-induced basis field (this comes next lecture). Another student remarks that in the box ``it should be $T_pU$ all over the place'';
Schuller: this is right, ``constitutes a basis of $T_pU$ is certainly right, $T_pM$ would be wrong because not for every point $p$ of $M$''. The statement is for any point of $U$, and the point must lie in the chart's domain (``the chart must be a neighbourhood of the point''); since the charts cover $M$, this is done for every point of $M$ by using another chart \ts{5}{1}{1:04:22}\ts{5}{1}{1:04:42}\ts{5}{1}{1:05:03}\ts{5}{1}{1:05:24}\ts{5}{1}{1:07:55}\ts{5}{1}{1:08:17}\ts{5}{1}{1:08:38}.
\begin{coursenote}[$T_pU=T_pM$]
For an open $U\ni p$ the tangent spaces $T_pU$ and $T_pM$ can be identified (a velocity of a curve in $M$ through $p$ is the velocity of a curve in $U$ after restricting $f$); in that reading the basis statement is also a statement about $T_pM$ for every $p\in U$, which is how the proof is used. (Not said in class.)
\end{coursenote}

\section{Change of vector components under a change of chart}\label{sec:c05-change}

The time is almost up; Schuller asks the audience to bear with him for fifteen more minutes to cover this important topic \ts{5}{1}{1:08:58}. In the older literature on differential geometry one finds statements of the type: ``tensors transform under a change of coordinate systems as follows, a vector transforms as follows''. Most precisely it must say that \emph{the components} of a vector transform. The vector is an abstract object, ``the velocity of the bird at that point in time'':
no change in the thinking about the coordinate system will change the velocity. If one considers the components in a particular chart and changes to another chart for the same velocity, the components change. So: change of vector components under a change of chart, not change of vector; ``the vector or tangent vector does not change under change of chart simply because it is defined independently of chart''
\ts{5}{1}{1:09:20}\ts{5}{1}{1:09:45}\ts{5}{1}{1:10:05}\ts{5}{1}{1:10:27}\ts{5}{1}{1:10:49}.

\paragraph{Terminology.} If $X\in T_pM$ one knows two things. First, there is a curve $\gamma$ through $p$ with $X=v_{\gamma,p}$ (by the definition of $T_pM$); one will in the future not always make the curve explicit, since different curves give the same vector. Second, as $T_pM$ is a vector space, there exist real numbers $X^1,\dots,X^d$ such that
\[
  X=X^i\,\frac{\partial}{\partial x^i}\Big|_p
\]
if a chart is chosen, but one could also pick any basis, not coming from a chart (``you can say: I like these three vectors to be the basis''); the $X^i$ are the components with respect to the chosen basis. ``If you have chosen a chart anyway, why not: you already committed the sin''
\ts{5}{1}{1:11:10}\ts{5}{1}{1:11:30}\ts{5}{1}{1:11:51}\ts{5}{1}{1:12:11}\ts{5}{1}{1:12:32}\ts{5}{1}{1:12:53}.

Let $(U,x)$ and $(V,y)$ be overlapping charts, $p\in U\cap V$, and $X\in T_pM$. Then $X=X^i_{(x)}\,\partial/\partial x^i|_p=X^j_{(y)}\,\partial/\partial y^j|_p$, the components with a subscript for the chart. How are they related? The abstract object in the middle is invariant \ts{5}{1}{1:13:13}\ts{5}{1}{1:13:34}\ts{5}{1}{1:13:56}\ts{5}{1}{1:14:17}\ts{5}{1}{1:14:39}\ts{5}{1}{1:15:00}.

\begin{theorem}[Change of vector components]\label{thm:c05-changeV}
For $p\in U\cap V$ the basis vectors transform as
\[
  \frac{\partial}{\partial x^i}\Big|_p=\frac{\partial y^j}{\partial x^i}\Big|_p\frac{\partial}{\partial y^j}\Big|_p,
\]
and hence the components of a vector $X\in T_pM$ as
\[
  X^j_{(y)}=\frac{\partial y^j}{\partial x^i}\Big|_p\,X^i_{(x)}.
\]
\end{theorem}
\begin{proof}
Act on $f\in C^\infty(M)$ and translate the symbol: $\dfrac{\partial f}{\partial x^i}\Big|_p=\partial_i(f\circ x^{-1})(x(p))$. Insert the identity $y^{-1}\circ y$:
\[
  \partial_i\bigl((f\circ y^{-1})\circ(y\circ x^{-1})\bigr)\bigl(x(p)\bigr).
\]
Now $f\circ y^{-1}$ goes from (a portion of) $\R^d$ to $\R$ and $y\circ x^{-1}$ is the old acquaintance, the chart transition map from $\R^d$ to $\R^d$. The multi-dimensional chain rule gives the inner derivative $\partial_i(y\circ x^{-1})^j\bigl(x(p)\bigr)$ times the outer derivative $\partial_j(f\circ y^{-1})$, evaluated at $(y\circ x^{-1})(x(p))=y(p)$:
\[
  \frac{\partial f}{\partial x^i}\Big|_p=\partial_i\bigl(y^j\circ x^{-1}\bigr)\bigl(x(p)\bigr)\cdot\partial_j\bigl(f\circ y^{-1}\bigr)\bigl(y(p)\bigr)=\frac{\partial y^j}{\partial x^i}\Big|_p\,\frac{\partial f}{\partial y^j}\Big|_p,
\]
where the $j$-th component of $y$ is the $j$-th coordinate function, so the first factor is of the form of the symbol with $f=y^j$ (``this $j$ comes because $y$ lands in $\R^d$''). This holds for all $f$, so the first formula holds. Insert it in $X=X^i_{(x)}\partial/\partial x^i|_p$ (renaming a summation index):
\[
  X^i_{(x)}\frac{\partial y^j}{\partial x^i}\Big|_p\frac{\partial}{\partial y^j}\Big|_p=X^j_{(y)}\frac{\partial}{\partial y^j}\Big|_p,
\]
and since the $\partial/\partial y^j|_p$ form a basis, one can compare coefficients (``numbers times the basis is zero, so the numbers are zero'').
\ts{5}{1}{1:15:21}\ts{5}{1}{1:15:41}\ts{5}{1}{1:16:03}\ts{5}{1}{1:16:23}\ts{5}{1}{1:16:45}\ts{5}{1}{1:17:05}\ts{5}{1}{1:17:25}\ts{5}{1}{1:17:46}\ts{5}{1}{1:18:06}\ts{5}{1}{1:18:27}\ts{5}{1}{1:18:47}\ts{5}{1}{1:19:09}\ts{5}{1}{1:19:30}\ts{5}{1}{1:19:51}\ts{5}{1}{1:20:12}\ts{5}{1}{1:20:34}\ts{5}{1}{1:20:55}.
\end{proof}
This is the formula one finds in books: the very same abstract object, ``the same physical vector, if you wish'', has in the $y$ chart the components obtained by calculating the numbers $\partial y^j/\partial x^i|_p$ and contracting. The change of coordinates can be wildly non-linear, but one goes to one point and evaluates the (wildly non-linear) dependence of $y$ on $x$ there:
one just gets a scheme of numbers, ``one might say a matrix, but we do not want to use that terminology'', and at each point the components are acted on by a constant linear map which depends only on the change of chart. So at each point there is a vector space acted on by a linear map, although the global chart transformation is non-linear \ts{5}{1}{1:20:55}\ts{5}{1}{1:21:16}\ts{5}{1}{1:21:37}\ts{5}{1}{1:21:57}\ts{5}{1}{1:22:18}.

\begin{aside}[Special relativity and Minkowski ``vector space'']
Does this not contradict special relativity, where one talks about the Minkowski vector space of spacetime points $x^\mu$? ``Yes, it flies in the face of everything'': the positions of special relativity are not a vector space either (the flat spacetime is not a vector space). The formulas work because one artificially restricts oneself to changing coordinate systems by \emph{linear} transformations, further restricted to Lorentz transformations.
That is a good idea for the individual velocity vectors, which transform under Lorentz transformations in each point of the flat tangent space by a linear map; but it does not mean that one mixes the points of the space only by linear maps: ``I want to be able to do special relativity in polar coordinates, if I wish'', and one can: changing from Cartesian to polar coordinates induces linear maps in the velocity tangent spaces
(and if one demands that observer velocities stay observer velocities, one can think about what one has to require). It is an over-structuralisation to equip the \emph{position space} of special relativity with a vector space structure, because then one also transforms position vectors with Lorentz transformations; in the tangent spaces it makes a lot of sense \ts{5}{1}{1:22:18}\ts{5}{1}{1:22:38}\ts{5}{1}{1:22:59}\ts{5}{1}{1:23:20}\ts{5}{1}{1:23:41}\ts{5}{1}{1:24:01}.
\end{aside}

\section{Cotangent spaces and the gradient}\label{sec:c05-cotangent}

Now the previous lecture pays off \ts{5}{1}{1:24:21}\ts{5}{1}{1:24:42}. At a point $p$ the tangent space $T_pM$ is just a vector space with the operations defined above, so one can take its dual.

\begin{definition}[New concepts: cotangent space, gradient]\label{def:c05-cotangent}
The \textbf{cotangent space} at $p$ is the dual vector space $T^*_pM:=(T_pM)^*$, the vector space of all linear maps $T_pM\to\R$ (\cref{def:c03-dual}). Its elements are \textbf{covectors} (or cotangent vectors).
For a smooth function $f\in C^\infty(M)$ define
\[
  \mathrm df_p\colon T_pM\to\R,\qquad\mathrm df_p(X):=X(f),
\]
the action of the tangent vector $X$ on the function $f$. This is linear in $X$ (by the definition of the vector space structure of \cref{thm:c05-vs}), so $\mathrm df_p\in T^*_pM$. It is called the \textbf{gradient of $f$ at $p$}.
\end{definition}
\ts{5}{1}{1:24:42}\ts{5}{1}{1:25:04}\ts{5}{1}{1:25:24}\ts{5}{1}{1:25:44}\ts{5}{1}{1:26:06}\ts{5}{1}{1:26:26}\ts{5}{1}{1:26:49}.
Why this complicated notation for something known, the directional derivative? Because the action of $X$ on $f$ is written as a linear map from the space of the $X$'s to $\R$, and then $\mathrm df_p$ is an element of the cotangent space. Notice that \emph{one does not need any chart to define the gradient}. The word ``cotangent'': the dual space of the tangent space; one writes the star up, as tradition [name explained only like this in the lecture].
(``Tangent'' and ``co-tangent'': the prefix \emph{co} reflects that covectors transform in the opposite way to vectors, see below [not said in class].)

\paragraph{Components of the gradient.} The gradient is a $(0,1)$-tensor over the vector space $T_pM$ (in the terminology of last time), a covector. How does one get the components of a $(0,1)$-tensor? Plug in the basis vectors with respect to which the components are wanted: here the chart-induced basis of $T_pM$ for a chart $(U,x)$. The result is
\[
  \mathrm df_p\Bigl(\frac{\partial}{\partial x^j}\Big|_p\Bigr)=\frac{\partial}{\partial x^j}\Big|_p(f)=\frac{\partial f}{\partial x^j}\Big|_p=\partial_j\bigl(f\circ x^{-1}\bigr)\bigl(x(p)\bigr),
\]
which one can compute by undergraduate analysis; one gets a number for each $j$, the components of the gradient. ``The gradient of a function is not a vector, the gradient of a function is a covector: there is nothing to be done about it''; for instance, under a change of chart its components transform differently from those of a vector \ts{5}{1}{1:27:10}\ts{5}{1}{1:27:32}\ts{5}{1}{1:27:52}\ts{5}{1}{1:28:13}\ts{5}{1}{1:28:33}\ts{5}{1}{1:28:54}\ts{5}{1}{1:29:14}.

\begin{theorem}[The dual basis of the chart-induced basis]\label{thm:c05-dualbasis}
Let $(U,x)$ be a chart and $p\in U$. The coordinate functions $x^a\colon U\to\R$ are smooth, and the gradients
\[
  \mathrm dx^1_p,\dots,\mathrm dx^d_p
\]
form a basis of $T^*_pM$, namely the \emph{dual basis} (\cref{def:c03-dualbasis}) of the chart-induced basis $\partial/\partial x^b|_p$ of $T_pM$:
\[
  \mathrm dx^a_p\Bigl(\frac{\partial}{\partial x^b}\Big|_p\Bigr)=\delta^a_b.
\]
\end{theorem}
\begin{proof}
By the definition of the gradient, $\mathrm dx^a_p\bigl(\partial/\partial x^b|_p\bigr)=\partial/\partial x^b|_p(x^a)=\partial_b(x^a\circ x^{-1})(x(p))=\delta^a_b$, as in the proof of \cref{thm:c05-basis}. A set of $d$ covectors in the $d$-dimensional space $T^*_pM$ which has this property is the dual basis of the given basis.
\ts{5}{1}{1:30:21}\ts{5}{1}{1:30:43}\ts{5}{1}{1:31:04}\ts{5}{1}{1:31:24}\ts{5}{1}{1:31:46}\ts{5}{1}{1:32:06}\ts{5}{1}{1:32:26}\ts{5}{1}{1:32:47}.
\end{proof}
(Schuller observes that in the dual space one could choose a basis totally independent of the one chosen in the tangent space, but constructing it as the gradients of the coordinate maps gives the dual basis.) This is the chart-induced basis for the cotangent space, given by the gradients of the coordinate maps.

\paragraph{Change of components of a covector.} A covector $\omega\in T^*_pM$ is an abstract object with ``physical reality'' and may be expanded as $\omega=\omega_{(x)\,i}\,\mathrm dx^i_p=\omega_{(y)\,j}\,\mathrm dy^j_p$. Going through the same motions as for vectors, with the multi-dimensional chain rule, one finds
\[
  \omega_{(y)\,i}=\frac{\partial x^j}{\partial y^i}\Big|_p\,\omega_{(x)\,j}.
\]
(Schuller does not prove it: ``you believe it works exactly like for the vectors''.) A mnemonic: the free index $i$ is down on the left, and the matrix has $y$ down and $x$ up, so that in the vector rule one has $x$ down/$y$ up; one has to remember which object goes where. The matrix for covector components ($x$ up, $y$ down) and the one for vector components ($y$ up, $x$ down) are \emph{inverses} of each other: this is on the problem sheet of the day. So the components of a covector transform \emph{the other way}, linearly, than those of a vector under the same change of chart; and again it is the components that change, the covector remains the same
\ts{5}{1}{1:29:37}\ts{5}{1}{1:30:00}\ts{5}{1}{1:33:09}\ts{5}{1}{1:33:30}\ts{5}{1}{1:33:51}\ts{5}{1}{1:34:12}\ts{5}{1}{1:34:33}\ts{5}{1}{1:34:53}\ts{5}{1}{1:35:14}.
\begin{coursenote}[The formula as Schuller writes it]
On the board the covector rule is given with $\partial x^j/\partial y^i$ (``$x$ up, $y$ down''); the version above uses the chain rule from $\mathrm dx^j_p=\partial x^j/\partial y^i\,\mathrm dy^i_p$, which is what the tutorial recalls (exercise 2 there). (Reconstructed from the description in the lecture, not written out in the transcript [?].)
\end{coursenote}

\begin{remark}[The gradient is not a vector]\label{rem:c05-gradient}
If one calls the gradient the ``gradient vector'' and believes the terminology, one contradicts the chain rule: write down the gradient, define its components as the partial derivatives, do the chain rule, which is undergraduate analysis, and one finds the wrong transformation behaviour, or at least wrong compared to the assumption that it transforms like a vector. ``It is not deep, it is not a special new mathematics that makes the gradient a covector: it is the chain rule of standard analysis that makes the gradient a covector and not a vector''
\ts{5}{1}{1:35:34}\ts{5}{1}{1:35:55}\ts{5}{1}{1:36:16}.
\end{remark}
The first half of the afternoon tutorial investigates this further; the lecture ends here, and the next one starts at ten past \ts{5}{1}{1:36:16}.

% ==================================================================
\section{Tutorial 5}\label{sec:t05-1}

Tutorial 5, 1:32:26, \seg{5}{2}{video}, by a tutor, not Schuller. The sheet (Sheet 5) is paraphrased, not copied; see \url{https://tales.mbivert.com/on-heraeus-winter-school-tutorials/} (and \url{https://tales.mbivert.com/heraeus/Sheet_05.pdf}). The first exercise is fifteen true/false statements in three groups of five \ts{5}{2}{0:00}\ts{5}{2}{0:20}.

\subsection*{Exercise 1: true or false}

\paragraph{(A) Tangent vectors.}
\begin{itemize}
  \item A tangent vector has a length and a direction: \textbf{false}; ``direction'' could be said, but a length is only defined once there is a metric, which in general there is not (a later lecture) \ts{5}{2}{0:20}\ts{5}{2}{0:40}\ts{5}{2}{1:02}.
  \item A tangent vector must not be zero: \textbf{false}; a curve that stays at $p$ for a while has $(f\circ\gamma)'=0$, so the velocity is the zero map $C^\infty(M)\to\R$ \ts{5}{2}{1:22}\ts{5}{2}{1:43}\ts{5}{2}{2:04}\ts{5}{2}{2:24}.
  \item A tangent vector maps a function on the manifold to a real number: \textbf{true}, this is the definition \ts{5}{2}{2:24}\ts{5}{2}{2:46}.
  \item A tangent vector is of the same dimension as the manifold: \textbf{false}; dimension is not defined for a vector, only for the whole vector space \ts{5}{2}{3:06}\ts{5}{2}{3:26}.
  \item A tangent vector arises as the velocity of some curve through the point: \textbf{true}, by the very definition of $T_pM$ \ts{5}{2}{3:26}\ts{5}{2}{3:48}.
\end{itemize}

\paragraph{(B) Tangent spaces.}
\begin{itemize}
  \item $T_pM$ has dimension $2\dim M$: \textbf{false}; it is $\dim M$, and $2\dim M$ is the dimension of the tangent bundle, which comes later \ts{5}{2}{4:09}\ts{5}{2}{4:30}.
  \item The tangent space can be defined at every point $p\in M$: \textbf{true}; the definition restricts nothing \ts{5}{2}{4:30}\ts{5}{2}{4:53}\ts{5}{2}{5:17}.
  \item $T_pM$ is no real vector space because its elements are only tangent vectors: \textbf{false}; the lecture equipped it with the operations and checked the axioms \ts{5}{2}{5:17}\ts{5}{2}{5:41}\ts{5}{2}{6:02}.
  \item $T_pM$ admits a linear bijection to $\R^d$ with componentwise operations: \textbf{true}; it is a $d$-dimensional real vector space and every such one is isomorphic to $\R^d$ \ts{5}{2}{6:02}\ts{5}{2}{6:22}\ts{5}{2}{6:44}\ts{5}{2}{7:07}.
  \item $T_pM$ has no tangent vector in common with $T_qM$ for $q\neq p$: \textbf{true}; the sets are defined by curves through $p$ and $q$ respectively \ts{5}{2}{7:07}\ts{5}{2}{7:28}.
\end{itemize}

\paragraph{(C) Charts and bases.}
\begin{itemize}
  \item The coordinate maps $x^i\colon U\to\R$ are only continuous, not differentiable: \textbf{false}. Down in the chart $x^i\circ x^{-1}$ is the $i$-th component of the identity of $\R^d$, which is differentiable, and since the chart transition maps are differentiable, so is it for any other chart: the notion on the manifold does not depend on the chart \ts{5}{2}{7:48}\ts{5}{2}{8:10}\ts{5}{2}{8:30}\ts{5}{2}{8:51}\ts{5}{2}{9:12}\ts{5}{2}{9:33}\ts{5}{2}{9:54}\ts{5}{2}{10:15}.
  \item The $\mathrm dx^i_p$ form a basis of $T^*_pM$: \textbf{true}; they are the dual basis of $\partial/\partial x^j|_p$, since $\mathrm dx^i_p(\partial/\partial x^j|_p)=\partial_j(x^i\circ x^{-1})(x(p))=\delta^i_j$ \ts{5}{2}{10:36}\ts{5}{2}{10:57}\ts{5}{2}{11:18}\ts{5}{2}{11:40}\ts{5}{2}{12:04}\ts{5}{2}{12:26}\ts{5}{2}{12:48}\ts{5}{2}{13:10}.
  \item There is no dual basis in the dual space of the basis $\partial/\partial x^i|_p$: \textbf{false}, it has just been given \ts{5}{2}{13:30}\ts{5}{2}{13:51}.
  \item The components of $X$ with respect to the chart-induced basis are $\mathrm dx^i_p(X)$: \textbf{true}; by the calculation with the chain rule $v_{\gamma,p}=\dot\gamma^i_x(\lambda_0)\partial/\partial x^i|_p$ and then $\mathrm dx^j_p(v_{\gamma,p})=v_{\gamma,p}(x^j)=\dot\gamma^i_x\delta^j_i$, which extracts the $j$-th component \ts{5}{2}{13:51}\ts{5}{2}{14:12}\ts{5}{2}{14:36}\ts{5}{2}{14:58}\ts{5}{2}{15:19}\ts{5}{2}{15:39}\ts{5}{2}{16:03}\ts{5}{2}{16:25}\ts{5}{2}{16:54}\ts{5}{2}{17:15}\ts{5}{2}{17:35}\ts{5}{2}{17:57}\ts{5}{2}{18:19}\ts{5}{2}{18:39}\ts{5}{2}{19:01}.
  \item $\mathrm dx^a_p(\partial/\partial x^a|_p)$ equals the dimension of the manifold: \textbf{true}; with the summation convention it is $\delta^a_a=\sum_{a=1}^d1=d$ \ts{5}{2}{19:23}\ts{5}{2}{19:44}\ts{5}{2}{20:05}\ts{5}{2}{20:28}\ts{5}{2}{20:49}.
\end{itemize}

\subsection*{Exercise 2: virtuosity with $\partial/\partial x^i|_p$}

Getting familiar with the symbol and related notions \ts{5}{2}{20:49}\ts{5}{2}{21:11}.
\begin{enumerate}
  \item \emph{Definition of $\partial f/\partial x^i|_p$.} For a smooth $f$ and a chart $(U,x)$ it looks like the undergraduate partial derivative, but $f$ lives on $M$, which carries no differentiable structure by itself. So go down by the chart: $f\circ x^{-1}$ is a map $\R^d\to\R$, take its $i$-th partial derivative (the derivative with respect to the $i$-th entry) and evaluate in the chart at $x(p)$:
        \[
          \frac{\partial f}{\partial x^i}\Big|_p:=\partial_i(f\circ x^{-1})(x(p))
        \]
        \ts{5}{2}{21:32}\ts{5}{2}{21:54}\ts{5}{2}{22:16}\ts{5}{2}{22:37}\ts{5}{2}{22:58}.
  \item \emph{Show that for overlapping charts $(U,x)$, $(V,y)$ one has $\dfrac{\partial x^a}{\partial y^m}\Big|_p\dfrac{\partial y^m}{\partial x^b}\Big|_p=\delta^a_b$.} Start with $\delta^a_b=\partial_b(x^a\circ x^{-1})(x(p))=\partial/\partial x^b|_p(x^a)$ (the previous exercise), insert $y^{-1}\circ y$, bracket $(x^a\circ y^{-1})\circ(y\circ x^{-1})$ and use the chain rule: the inner derivative of the $m$-th component gives
        $\partial_b(y^m\circ x^{-1})(x(p))=\partial y^m/\partial x^b|_p$ and the outer one gives $\partial_m(x^a\circ y^{-1})(y(p))=\partial x^a/\partial y^m|_p$, where $x^{-1}$ cancels $x$ so that one evaluates at $y(p)$. Picking the $m$-th component at the end of the inner function or ``in front'' makes no difference
        \ts{5}{2}{22:58}\ts{5}{2}{23:20}\ts{5}{2}{23:42}\ts{5}{2}{24:06}\ts{5}{2}{24:27}\ts{5}{2}{24:49}\ts{5}{2}{25:13}\ts{5}{2}{25:33}\ts{5}{2}{25:54}\ts{5}{2}{26:14}\ts{5}{2}{26:35}\ts{5}{2}{26:56}\ts{5}{2}{27:17}\ts{5}{2}{27:38}\ts{5}{2}{27:59}\ts{5}{2}{28:21}\ts{5}{2}{28:43}\ts{5}{2}{29:04}.
        This is the statement that the two matrices of \cref{thm:c05-changeV} and of the covector rule are inverse to each other (the problem sheet announced in the lecture).
  \item \emph{The chain rule $\dfrac{\partial f}{\partial x^i}\Big|_p=\dfrac{\partial y^m}{\partial x^i}\Big|_p\dfrac{\partial f}{\partial y^m}\Big|_p$.} Translate the left side to $\partial_i(f\circ x^{-1})(x(p))$, insert the identity $y^{-1}\circ y$ at the right place, bracket and use the chain rule; this is the proof of \cref{thm:c05-changeV} \ts{5}{2}{29:28}\ts{5}{2}{29:49}\ts{5}{2}{30:09}\ts{5}{2}{30:31}\ts{5}{2}{30:56}\ts{5}{2}{31:21}\ts{5}{2}{31:46}\ts{5}{2}{32:12}\ts{5}{2}{32:33}\ts{5}{2}{32:53}.
  \item \emph{Do the $d$ numbers $\partial f/\partial x^i|_p$ form the components of a tensor, and of which valence and rank?} Yes, of a $(0,1)$-tensor (a covector), of rank $0+1=1$: expand a covector $\omega=\omega_{(x)\,i}\,\mathrm dx^i_p=\omega_{(y)\,j}\,\mathrm dy^j_p$, the components must satisfy $\omega_{(x)\,i}=\dfrac{\partial y^j}{\partial x^i}\Big|_p\omega_{(y)\,j}$, and
        by the previous part the numbers $\partial f/\partial x^i|_p$ obey exactly this. This is the gradient $\mathrm df_p$ of the lecture \ts{5}{2}{32:53}\ts{5}{2}{33:17}\ts{5}{2}{33:39}\ts{5}{2}{33:59}\ts{5}{2}{34:20}\ts{5}{2}{34:41}\ts{5}{2}{35:01}\ts{5}{2}{35:23}\ts{5}{2}{35:43}\ts{5}{2}{36:04}.
\end{enumerate}

\subsection*{Exercise 3: transformation of vector components, the pedestrian way}

Real components of real vectors \ts{5}{2}{36:04}\ts{5}{2}{36:26}. On $\R^2$ with the standard topology take the atlas of two charts: $x=\mathrm{id}$ and $y(a,b)=(a,\,b+a^3)$ [the sheet's map is read from the computation of $y^{-1}$ below, reconstructed [?]]. Let $\gamma(\lambda)=(\lambda,-\lambda)$ [the curve, read from the computation of $x\circ\gamma$ below]
\ts{5}{2}{36:47}\ts{5}{2}{37:07}.
\begin{enumerate}
  \item \emph{The four numbers $\partial x^i/\partial y^j|_p$.} By definition $\partial x^i/\partial y^j|_p=\partial_j(x^i\circ y^{-1})(y(p))$. Invert $y$: $y^1=a$, $y^2=b+a^3$, so $y^{-1}(u,v)=(u,\,v-u^3)$, and $x\circ y^{-1}=y^{-1}$. For $p=(a,b)$ one gets
        \[
          \frac{\partial x^1}{\partial y^1}\Big|_p=1,\qquad\frac{\partial x^2}{\partial y^1}\Big|_p=-3a^2,\qquad\frac{\partial x^1}{\partial y^2}\Big|_p=0,\qquad\frac{\partial x^2}{\partial y^2}\Big|_p=1.
        \]
        The tutor writes $-2a^2$ first and corrects the slip only in part 3: the derivative of $-u^3$ is $-3u^2$ (the corrected value is the one used below) \ts{5}{2}{37:28}\ts{5}{2}{37:49}\ts{5}{2}{38:10}\ts{5}{2}{38:31}\ts{5}{2}{38:52}\ts{5}{2}{39:13}\ts{5}{2}{39:41}\ts{5}{2}{40:03}\ts{5}{2}{40:24}\ts{5}{2}{40:45}\ts{5}{2}{41:08}\ts{5}{2}{41:29}\ts{5}{2}{41:50}\ts{5}{2}{42:17}\ts{5}{2}{42:38}\ts{5}{2}{43:00}\ts{5}{2}{43:21}.
  \item \emph{Components of $v_{\gamma,p}$ in the two charts.} With $\dot\gamma^i_x(\lambda_0)=(x^i\circ\gamma)'(\lambda_0)$: since $x$ is the identity, $x\circ\gamma(\lambda)=(\lambda,-\lambda)$ and the components are $\dot\gamma^1_x=1$, $\dot\gamma^2_x=-1$. In the chart $y$: $y\circ\gamma(\lambda)=(\lambda,-\lambda+\lambda^3)$, so $\dot\gamma^1_y=1$ and $\dot\gamma^2_y=-1+3\lambda_0^2$
        \ts{5}{2}{43:51}\ts{5}{2}{44:12}\ts{5}{2}{44:33}\ts{5}{2}{44:56}\ts{5}{2}{45:17}\ts{5}{2}{45:40}\ts{5}{2}{46:00}\ts{5}{2}{46:21}\ts{5}{2}{46:42}\ts{5}{2}{47:05}\ts{5}{2}{47:26}\ts{5}{2}{47:47}\ts{5}{2}{48:08}.
  \item \emph{Obtain $\dot\gamma_x$ from $\dot\gamma_y$.} Insert $y^{-1}\circ y$ in $(x^i\circ\gamma)'(\lambda_0)$ and use the chain rule: $\dot\gamma^i_x(\lambda_0)=\dfrac{\partial x^i}{\partial y^m}\Big|_{\gamma(\lambda_0)}\dot\gamma^m_y(\lambda_0)$, the transformation law of vector components. Check: for the first component $1\cdot1+0\cdot(-1+3\lambda_0^2)=1$, as before. For the second, with $a=\lambda_0$:
        $-3\lambda_0^2\cdot1+1\cdot(-1+3\lambda_0^2)=-1$, again the same as the direct computation ($-1$). The tutor first plugs in the wrong $-2a^2$ and finds a mismatch, spots the mistake of the first part and corrects it: ``now the right answer''
        \ts{5}{2}{48:28}\ts{5}{2}{48:49}\ts{5}{2}{49:12}\ts{5}{2}{49:33}\ts{5}{2}{49:54}\ts{5}{2}{50:15}\ts{5}{2}{50:37}\ts{5}{2}{51:04}\ts{5}{2}{51:29}\ts{5}{2}{51:51}\ts{5}{2}{52:14}\ts{5}{2}{52:35}\ts{5}{2}{52:56}\ts{5}{2}{53:19}\ts{5}{2}{53:44}\ts{5}{2}{54:06}\ts{5}{2}{54:27}\ts{5}{2}{54:49}\ts{5}{2}{55:09}\ts{5}{2}{55:32}\ts{5}{2}{55:53}\ts{5}{2}{56:14}\ts{5}{2}{56:35}\ts{5}{2}{56:57}\ts{5}{2}{57:17}\ts{5}{2}{57:40}\ts{5}{2}{58:01}.
\end{enumerate}
\begin{coursenote}[Slip in the tutorial]
In the first part the tutor reads off $-2a^2$ for $\partial x^2/\partial y^1|_p$; the correct value is $-3a^2$ (derivative of $v-u^3$ in $u$). The error is noticed during the check at \ts{5}{2}{56:57}. The lecture-notes values above are the corrected ones.
\end{coursenote}

\subsection*{Exercise 4: the gradient annihilates curves in a level set}

The \emph{level set} of $f\colon M\to\R$ for a constant $c$ is $\{p\in M\mid f(p)=c\}$ (picture: level lines $c_1,c_2,c_3$) \ts{5}{2}{58:23}\ts{5}{2}{58:43}\ts{5}{2}{59:04}\ts{5}{2}{59:25}.
\begin{enumerate}
  \item A curve $\gamma$ lies in the level set iff $f(\gamma(\lambda))=c$ for all $\lambda$, i.e.\ $f\circ\gamma$ is the constant function $c$ \ts{5}{2}{59:45}\ts{5}{2}{1:00:06}\ts{5}{2}{1:00:27}\ts{5}{2}{1:00:47}\ts{5}{2}{1:01:08}.
  \item Then for a point $p$ of the level set and $\gamma(\lambda_0)=p$: $\mathrm df_p(v_{\gamma,p})=v_{\gamma,p}(f)=(f\circ\gamma)'(\lambda_0)=0$, the derivative of a constant. So the gradient annihilates the velocities of curves which stay in a level set
        \ts{5}{2}{1:01:34}\ts{5}{2}{1:01:58}\ts{5}{2}{1:02:18}\ts{5}{2}{1:02:42}\ts{5}{2}{1:03:04}\ts{5}{2}{1:03:24}.
\end{enumerate}

\subsection*{Exercise 5: is there a well-defined sum of curves?}

The sheet follows the construction of the lecture to see whether it defines a \emph{sum of curves} and not only of velocities \ts{5}{2}{1:03:47}\ts{5}{2}{1:04:07}. The exercise is long (about 30 minutes, five parts); its core is kept here and the full solution is moved to the appendix (\cref{app:c05-curves}).
The atlas is $\R^2$ with $x=\mathrm{id}$ and $y(a,b)=(a,be^a)$, a smooth atlas (the transition maps $(a,b)\mapsto(a,be^a)$ and $(a,b)\mapsto(a,be^{-a})$ are smooth). The curves $\gamma(\lambda)=(\lambda,1)$ and $\delta(\lambda)=(1,\lambda)$ meet at $p=(1,1)$ for $\lambda_0=\lambda_1=1$.
Without a chart there is no sum $\gamma+\delta$, since $\R^2$ here is the manifold and not a vector space. With the construction of the lecture, in the chart $x$ one gets $\sigma_x(\lambda)=(1+\lambda,1+\lambda)$, in the chart $y$ one gets $\sigma_y(\lambda)=(1+\lambda,1+\lambda e^{-\lambda})$: \emph{different curves}, so the sum of curves is chart dependent and useless. But the components of their velocities in the same chart $x$ are
$(1,1)$ for both: $\sigma_x'(0)=(1,1)$ and $\sigma_y'(0)=(1,\,e^{-\lambda}(1-\lambda)|_0)=(1,1)$; hence the two velocities coincide, as the lecture says \ts{5}{2}{1:04:07}\ts{5}{2}{1:26:57}\ts{5}{2}{1:31:55}.

% ==================================================================
\section*{Key concepts}
\begin{keyconcepts}
\textbf{Retrieval questions.}
\begin{enumerate}
  \item Define the velocity of a curve at a point as a map on $C^\infty(M)$. What is its picture in terms of a temperature and a child?
  \item Why does the intrinsic definition of $T_pM$ not need an ambient space? What do the embedding theorems say?
  \item How does one add two velocities at $p$ and scale one? Why is a chart needed for the sum, why is the result chart independent, and why is the sum of the curves itself not?
  \item Define the symbol $\partial/\partial x^i|_p$. What is it a map on, and what must one always do when one meets it?
  \item State and prove that the chart-induced vectors form a basis of $T_pM$. What is the dimension of $T_pM$?
  \item State the transformation law of vector components under a change of chart. Does the vector change? Why is it not a contradiction with special relativity?
  \item Define the cotangent space and the gradient $\mathrm df_p$. Why does it need no chart? Why is the gradient a covector and not a vector?
  \item What is the dual basis of $\partial/\partial x^i|_p$, and how do the components of a covector transform?
\end{enumerate}
\textbf{Concepts.}
\begin{itemize}
  \item Velocity $v_{\gamma,p}(f)=(f\circ\gamma)'(\lambda_0)$ (\cref{def:c05-velocity}); tangent space $T_pM$ (\cref{def:c05-tangent}) and its vector space structure (\cref{thm:c05-vs}).
  \item Components $\dot\gamma^i_x(0)$ and the chart-induced basis $\partial/\partial x^i|_p$ (\cref{def:c05-chartbasis}, \cref{def:c05-inducedbasis}, \cref{thm:c05-basis}); $\dim T_pM=\dim M$.
  \item Change of components $X^j_{(y)}=\dfrac{\partial y^j}{\partial x^i}\Big|_pX^i_{(x)}$ (\cref{thm:c05-changeV}).
  \item Cotangent space $T^*_pM$, gradient $\mathrm df_p(X)=X(f)$ (\cref{def:c05-cotangent}), dual basis $\mathrm dx^a_p$ (\cref{thm:c05-dualbasis}), covector components transform with the inverse matrix (\cref{rem:c05-gradient}).
  \item The gradient annihilates velocities of curves in a level set; the sum of curves is chart dependent, the sum of velocities is not (\cref{sec:t05-1}).
\end{itemize}
\end{keyconcepts}

\flushside
\end{paracol}
